% =====================================================================
%  Thème des diapositives « Introduction à AWS »
%  Support de formation : un module par séance, sommaire rappelé à chaque
%  module, diapos de TP et de questions, numérotation « module - page ».
%  Figures : ../static/figures (TikZ + icônes officielles AWS).
% =====================================================================
\usepackage{fontspec}
\setsansfont{Inter}[Scale=0.95]
\setmonofont{DejaVu Sans Mono}[Scale=0.78]
\usefonttheme{professionalfonts}
\usepackage[french]{babel}
\usepackage{tikz}
\usetikzlibrary{positioning,calc,arrows.meta,fit}
\usepackage{tcolorbox}
\usepackage{booktabs}
\usepackage{tabularx}
\usepackage{array}
\usepackage{colortbl}
\usepackage{listings}
\usepackage{fontawesome5}
\graphicspath{{../static/figures/}{../figures/icones/}}
\hyphenpenalty=10000
\exhyphenpenalty=10000

% ---- couleurs --------------------------------------------------------
\definecolor{Squid}{HTML}{232F3E}
\definecolor{Orange}{HTML}{EC7211}
\definecolor{Encre}{HTML}{16191F}
\definecolor{Gris}{HTML}{5F6B7A}
\definecolor{Ligne}{HTML}{D1D5DB}
\definecolor{Doux}{HTML}{F2F3F3}
\definecolor{Lien}{HTML}{0972D3}
\definecolor{Rouge}{HTML}{D91515}
\definecolor{Vert}{HTML}{037F0C}

\hypersetup{colorlinks=true, urlcolor=Lien, linkcolor=Squid}
\setbeamercolor{normal text}{fg=Encre}
\setbeamercolor{structure}{fg=Squid}
\setbeamercolor{alerted text}{fg=Orange}
\setbeamercolor{item}{fg=Squid}
\setbeamertemplate{navigation symbols}{}
\setbeamertemplate{itemize item}{\raisebox{.18em}{\color{Squid}\rule{.36em}{.36em}}}
\setbeamertemplate{itemize subitem}{\raisebox{.22em}{\color{Gris}\rule{.26em}{.26em}}}
\setbeamertemplate{enumerate item}{\color{Squid}\bfseries\insertenumlabel.}
\setlength{\leftmargini}{1.1em}
\setlength{\leftmarginii}{1.1em}
\setbeamerfont{itemize/enumerate subbody}{size=\normalsize}

% ---- numérotation « module - page » ------------------------------------
\newcount\debutmodule \debutmodule=0
\newcommand{\numdiapo}{%
  \ifnum\value{section}>0 \thesection\,-\,\the\numexpr\value{framenumber}-\debutmodule\relax
  \else\insertframenumber\fi}

% ---- titre de diapo : capitales, bandeau orange, module en haut à droite --
\setbeamertemplate{frametitle}{%
  \vspace*{.7em}%
  {\usebeamerfont{frametitle}\Large\bfseries\color{Squid}\MakeUppercase{\insertframetitle}}\par
  \vspace*{-.1em}}
\setbeamertemplate{headline}{%
  \begin{tikzpicture}[remember picture, overlay]
    \fill[Orange] (current page.north west) rectangle ([yshift=-.11cm]current page.north east);
  \end{tikzpicture}}
\setbeamertemplate{footline}{%
  \begin{tikzpicture}[remember picture, overlay]
    \node[anchor=west, font=\scriptsize, text=Gris] at ([xshift=.6cm,yshift=.35cm]current page.south west) {\textbf{\numdiapo}\quad Introduction à AWS};
    \ifnum\value{section}>0
      \node[anchor=east, font=\scriptsize, text=Gris] at ([xshift=-.6cm,yshift=.35cm]current page.south east) {Module \thesection\ · \insertsectionhead};
    \fi
  \end{tikzpicture}}

% ---- source et lien vers la documentation ----------------------------------
\newcommand{\source}[1]{%
  \begin{tikzpicture}[remember picture, overlay]
    \node[anchor=south west, font=\tiny, text=Gris, text width=13.6cm, align=left, inner sep=0pt]
      at ([xshift=.6cm,yshift=.72cm]current page.south west) {Source : #1};
  \end{tikzpicture}}
\newtcolorbox{encadrelien}{colback=Doux, colframe=Doux, sharp corners, boxrule=0pt,
  left=6pt, right=6pt, top=3pt, bottom=3pt, fontupper=\small}
\newcommand{\lien}[1]{\begin{encadrelien}\url{#1}\end{encadrelien}}

% ---- mise en garde en ligne --------------------------------------------------
\newcommand{\attention}{{\color{Orange}\faExclamationTriangle}\ }

% ---- terminal pour une commande ---------------------------------------------
\newtcolorbox{terminal}{colback=Squid, colframe=Squid, sharp corners, boxrule=0pt,
  left=8pt, right=8pt, top=4pt, bottom=4pt, fontupper=\ttfamily\small\color{white}}

% ---- JSON ---------------------------------------------------------------------
\lstdefinelanguage{json}{morestring=[b]", stringstyle=\color{Lien}, literate={:}{{\color{Gris}:}}1}
\lstset{basicstyle=\ttfamily\footnotesize, columns=fullflexible, keepspaces=true,
  backgroundcolor=\color{Doux}, frame=none, xleftmargin=6pt, xrightmargin=6pt,
  aboveskip=4pt, belowskip=4pt, showstringspaces=false, upquote=true,
  literate={é}{{\'e}}1 {è}{{\`e}}1 {à}{{\`a}}1}

% ---- flèche des schémas propres aux diapos ---------------------------------------
\tikzset{fl/.style={-{Stealth[length=2.4mm]}, line width=1pt, draw=#1}, fl/.default=Encre}

% ---- liste des modules (sommaire) ------------------------------------------------
\newcommand{\listemodules}{%
  {Introduction à AWS},
  {Sécurité et gestion des accès},
  {Calcul : Amazon EC2},
  {Stockage, bases de données et messages},
  {Projet final}}
\newcommand{\sommaire}[1]{%
  \begin{frame}{Sommaire}
    \begin{tikzpicture}[x=1cm,y=1cm]
      \foreach \titre [count=\n] in \listemodules {
        \ifnum\n=#1
          \node[anchor=west, font=\large\bfseries, text=Orange] at (0,-\n*.82) {\n.\enspace\titre};
          \fill[Orange] (-.35,-\n*.82-.2) rectangle ++(.1,.4);
        \else
          \node[anchor=west, font=\large, text=Gris] at (0,-\n*.82) {\n.\enspace\titre};
        \fi
      }
    \end{tikzpicture}
  \end{frame}}

% ---- page de module (à chaque \section) --------------------------------------------
\AtBeginSection{%
  \global\debutmodule=\numexpr\value{framenumber}\relax
  {\setbeamertemplate{headline}{\begin{tikzpicture}[remember picture, overlay]\fill[Orange] (current page.north west) rectangle ([yshift=-.11cm]current page.north east);\end{tikzpicture}}
   \begin{frame}[plain]
     \begin{tikzpicture}[remember picture, overlay]
       \fill[Orange] (current page.north west) rectangle ([yshift=-.11cm]current page.north east);
       \node[circle, fill=Squid, text=white, font=\fontsize{30}{30}\selectfont\bfseries, minimum size=2.2cm] at ([yshift=1.1cm]current page.center) {\thesection};
       \node[font=\fontsize{24}{28}\selectfont\bfseries, text=Squid, text width=13cm, align=center] at ([yshift=-.9cm]current page.center) {\MakeUppercase{\insertsectionhead}};
     \end{tikzpicture}
   \end{frame}}
  \sommaire{\value{section}}}

% ---- diapo de TP -----------------------------------------------------------------
\newcommand{\diapotp}[3]{% numéro, titre, éléments \item
  \begin{frame}[plain]
    \begin{tikzpicture}[remember picture, overlay]
      \fill[Orange] (current page.north west) rectangle ([yshift=-.11cm]current page.north east);
      \fill[Squid] (current page.north west) rectangle ([xshift=5.2cm]current page.south west);
      \node[font=\fontsize{40}{40}\selectfont\bfseries, text=white] at ([xshift=2.6cm,yshift=.5cm]current page.west) {TP #1};
      \node[font=\large, text=Ligne, text width=4.2cm, align=center] at ([xshift=2.6cm,yshift=-.7cm]current page.west) {#2};
      \node[anchor=west, text width=8.6cm] at ([xshift=6.1cm]current page.west) {%
        \begin{minipage}{8.6cm}\begin{itemize}\setlength{\itemsep}{.6em}#3\end{itemize}\end{minipage}};
    \end{tikzpicture}
  \end{frame}}

% ---- icône de service avec légende -------------------------------------------------------
\newcommand{\service}[3][1.2cm]{%
  \begin{tabular}{@{}c@{}}\includegraphics[width=#1]{#2}\\[.15em]\parbox[t]{2.7cm}{\centering\small #3}\end{tabular}}
